\section{Method}
\subsection{Multiband ECoG Representation}
Let $x_c[n]$ denote the 1,000-Hz ECoG signal at electrode $c$. For each independently processed recording segment, we standardize each electrode by that segment's mean and standard deviation, then subtract the pointwise median across all available electrodes. Segment-level standardization uses unlabeled ECoG from the complete segment, including a complete test segment at inference. All subsequent supervised electrode selection and feature/target scaling use training data.

The high-frequency branch applies a 40--300-Hz FIR bandpass and notches at 60-Hz harmonics below 500~Hz. We compute seven-cycle Morlet power at 40 logarithmically spaced centers,
\begin{equation}
f_k=40\left(\frac{300}{40}\right)^{k/39},\qquad k=0,\ldots,39,
\end{equation}
with FFT convolution, non-zero-mean wavelets, and decimation by ten. The low-frequency branch operates on the referenced signal before the high-frequency bandpass. Fourth-order, zero-phase Butterworth filters extract $[0.5,4]$, $[4,8]$, $[8,13]$, and $[13,30]$~Hz waveforms. Their signed outputs, also sampled every ten points, retain polarity instead of converting voltage to power.

The resulting 100-Hz representation contains 44 features per electrode. Sixteen training-derived electrodes are retained. For electrode selection, we rank each electrode by maximum absolute same-time feature/target correlation over the four scored fingers, sampling the training interior every five points; final BCI runs retain their fixed training-derived sets. Features are centered and scaled independently per electrode and feature using a training-fitted RobustScaler with percentile range $(0.1,0.9)$ and unit-variance correction. These settings denote the 0.1st and 0.9th percentiles. Signed low-frequency features are multiplied by 200 after scaling; power features have unit gain. Fitting excludes the first 225 and last 25 points of the 100-Hz training segment to reduce preprocessing-edge influence. Test features reuse these fitted statistics.

Training targets are reconstructed from every 40th point of the stored 1,000-Hz glove signal, corresponding to its 25-Hz support. We append the final support value and use cubic interpolation to 100~Hz. Each target is then scaled using its training-interior minimum and range. Input and target windows use identical time indices; no explicit temporal offset is introduced.



\subsection{Temporal Encoder--Decoder with Recurrent Context}
The input $\mathbf X\in\mathbb R^{B\times16\times44\times T}$ is flattened across electrodes and features. A width-three convolution reduces 704 channels to 64. Five encoder stages use output widths $(64,128,256,512,512)$, kernels $(7,7,5,5,5)$, dilations $(1,2,3,1,2)$, and factor-two max pooling. Each convolution block has same padding, channel-wise layer normalization at each time point, GELU activation, and dropout of 0.1. The bottleneck is $\mathbf Z\in\mathbb R^{B\times512\times T/32}$.

We normalize the bottleneck channels and apply a single-layer BiGRU with 128 hidden units in each direction. A linear projection maps its 256-dimensional output back to 512 channels. The residual update is
\begin{equation}
\widetilde{\mathbf Z}=\mathbf Z+g\,\operatorname{Proj}
\bigl(\operatorname{BiGRU}(\operatorname{LN}(\mathbf Z))\bigr),
\end{equation}
where $\operatorname{Proj}$ is a learned affine projection, time/channel axes are transposed for recurrence, and the learned scalar $g$ starts at 0.1. This concentrates recurrent computation on the compressed time axis while preserving the convolutional residual path.

Five mirrored decoder stages use convolution blocks followed by width-two, stride-two transposed convolutions. At each scale, the upsampled output is concatenated with the corresponding encoder feature. A final $1\times1$ convolution maps 128 channels to five trajectories at the input resolution. The model has 6,808,198 trainable parameters; $T$ is divisible by 32. Figure~\ref{fig:model} summarizes the architecture.

\subsection{Learning and Continuous Reconstruction}
We train one seed-42 model per subject with random 512-point windows (5.12~s), batch size 64, and 8,192 sampled windows per epoch. For normalized prediction $\hat{\mathbf y}_{bf}$ and target $\mathbf y_{bf}$ along a window's time axis, the objective is
\begin{equation}
\begin{split}
\mathcal L={}&\frac{\|\hat{\mathbf Y}-\mathbf Y\|_F^2}{2\cdot5BT}\\
&+\frac{1}{2\cdot5B}\sum_{b,f}
\left(1-\frac{\langle\hat{\mathbf y}_{bf},\mathbf y_{bf}\rangle}
{\max(\|\hat{\mathbf y}_{bf}\|_2,\epsilon)\max(\|\mathbf y_{bf}\|_2,\epsilon)}\right).
\end{split}
\end{equation}
Here $\epsilon=10^{-8}$, as in the implemented cosine similarity. AdamW uses learning rate $8.42\times10^{-5}$, weight decay $10^{-4}$, and gradient-norm clipping at one. Training starts from random weights, with shared tensors initialized consistently across feature controls and the added alpha slot initialized afresh. No pretrained weights are used.

Inference uses consecutive, non-overlapping 2,048-point blocks (20.48~s). Every output is retained, the final incomplete block is reflection-padded and trimmed, and GRU state starts at zero in each block. Predictions are inverse-scaled to glove units and linearly interpolated to the stored 1,000-Hz axis; samples beyond the final 100-Hz support use its endpoint value. No ensemble, Gaussian smoothing, amplitude calibration, output floor, or clipping is applied.

Centered preprocessing, segment-wide normalization, and bidirectional recurrence use future ECoG; RecurFlex therefore performs offline reconstruction.
